% TOP_PROBLEM: {"id": 30006834, "title": "Resolvent Degree of the General Degree-Seven Polynomial (Algebraic Form of Hilbert's Thirteenth Problem)", "category_id": 4, "category": {"id": 4, "name": "algebra", "display_name": "Algebra", "description": "Group theory, ring theory, field theory, and algebraic structures.", "slug": "algebra", "order_index": 4, "created_at": "2026-07-31T15:26:25.671Z"}, "status": "open"}
% FIRST_POSTED: 2026-09-25
% ENTRY_KIND: statement_only
% !TeX program = lualatex
% Definition and sources only; this file is not an LLM solution attempt.
\documentclass[11pt]{article}
\usepackage[margin=25mm]{geometry}
\usepackage{fontspec}
\setmainfont{Latin Modern Roman}
\usepackage{amsmath,amssymb,mathtools}
\usepackage{unicode-math}
\setmathfont{Latin Modern Math}
\usepackage{xurl,hyperref}
\hypersetup{colorlinks=true,urlcolor=blue}
\setcounter{secnumdepth}{0}
\setlength{\parindent}{0pt}
\setlength{\parskip}{5pt}
\setlength{\emergencystretch}{3em}
\begin{document}
\section*{\texorpdfstring{Resolvent Degree of the General Degree-Seven Polynomial (Algebraic Form of Hilbert's Thirteenth Problem)}{Resolvent Degree of the General Degree-Seven Polynomial (Algebraic Form of Hilbert's Thirteenth Problem)}}\label{30006834}
\subsection{Definitions and mathematical statement}
Resolvent degree over ${\mathbb{C}}$ measures how many independent parameters are needed in the algebraic functions used to solve a generic equation. Formally, it is the least $d$ for which the generic root cover is dominated by a finite tower of covers whose steps arise by rational pullback from algebraic covers of varieties of dimension at most $d$, using the standard birational convention of the source. For the general polynomial
$X^7+a_1X^6+\cdots+a_7$ with algebraically independent coefficients, the target is
\[
 \operatorname{rd}_{{\mathbb{C}}}(7)=3,
\]
or, using the upper bound incorporated in the catalogue, the obstruction $\operatorname{rd}_{{\mathbb{C}}}(7)>2$. This is an algebraic-function problem, not the continuous-superposition version of Hilbert's thirteenth problem.
\par\smallskip\textit{Formulation and concept references:} \href{https://iopscience.iop.org/article/10.1070/RM2004v059n01ABEH000698}{[S1]}.

\subsection{Short English statement}
Does solving the generic degree-seven polynomial necessarily require algebraic functions of at least three independent variables, even when finite towers of such functions are allowed?
\subsection{Sources}
\begin{itemize}
\item[S1]A. G. Vitushkin, "On Hilbert's thirteenth problem and related questions", Russian Mathematical Surveys 59 (2004), no. 1; DOI 10.1070/RM2004v059n01ABEH000698.
\url{https://iopscience.iop.org/article/10.1070/RM2004v059n01ABEH000698}.
\textit{Formulation source.}
\end{itemize}
\end{document}
